% =====================================================================
% Introduction (CVPR style, 5 paragraphs, ~1 column)
% Structure:
%   P1. Motivation: why multi-camera person search matters in the wild.
%   P2. Gap: classical pipelines are single-shot, modality-fixed, and
%       opaque; clothing change, occlusion, and cross-camera transit
%       expose the limit of one-shot matching.
%   P3. Insight: person search is a multi-step tool-using decision
%       problem; LLM agents excel at exactly that.
%   P4. Approach: SearchCop = Planner / Executor / Critic over a tool
%       library + spatio-temporal prior + self-correcting loop, with
%       five enumerated contributions.
%   P5. Results headline + roadmap of the paper.
% =====================================================================
\section{Introduction}
\label{sec:intro}

Locating a person of interest across a network of surveillance cameras is a
foundational capability for public safety, missing-person investigation,
and large-scale forensic analysis. In contrast to image-level person
re-identification (Re-ID), the operational task analysts actually face is
\emph{multi-camera person search}: given a short query (an image, a video
clip, or even a textual description), determine \emph{when} and \emph{from
which camera} the same individual was seen, and return a ranked list of
candidate tracklets together with an explanation that a human can verify.
Modern monitoring deployments comprise hundreds of overlapping or sparsely
connected cameras, and queries arrive with heterogeneous quality: a frontal
crop here, a back-view tracklet there, sometimes only a witness sentence.
This setting demands a search system that is not only accurate, but also
\emph{adaptive}, \emph{cost-aware}, and \emph{auditable}.

The dominant approach today reduces multi-camera search to a single-shot
feature-matching procedure: encode the query with a Re-ID
network~\cite{he2021transreid,li2023clipreid} or a gait
network~\cite{chao2019gaitset,fan2023opengait}, retrieve top-$K$ candidates
from a fixed-modality FAISS index, and optionally late-fuse the two
modalities with hand-tuned weights. This recipe ignores three structural
properties of the problem. First, query quality is highly variable -- the
right encoder for a clean walking sequence is the wrong encoder for a
heavily occluded crop, yet existing systems commit to a modality at
indexing time. Second, multi-camera networks have rich spatio-temporal
priors: a person seen at camera~A at time $t$ cannot physically appear at
a non-adjacent camera~B within a few seconds, yet most retrieval pipelines
discard topology entirely~\cite{davila2023mevid,zheng2017prw}. Third,
clothing-change benchmarks~\cite{yang2019prcc,qian2020ltcc} reveal that any
single modality, however carefully fine-tuned, fails on a non-trivial slice
of queries; the fix is not a stronger encoder, but a \emph{decision
policy} that knows when to switch.

Our key insight is that multi-camera person search is more naturally cast
as a \emph{multi-step tool-using decision problem} than as a one-shot
retrieval problem. Recent advances in tool-using large language model
(LLM) agents -- from ReAct~\cite{yao2023react} and
Toolformer~\cite{schick2023toolformer} to Visual
Programming~\cite{gupta2023visprog} and Chameleon~\cite{lu2023chameleon}
-- show that LLMs, when given a structured tool library, can plan
non-trivial sequences of vision/database calls and produce traceable
rationales. We bring this paradigm to person search. Crucially, the same
encoders that prior work uses as monolithic pipelines (gait, Re-ID,
vision-language captioners) become \emph{callable tools} in the LLM's
action space, and the same camera-topology priors that are usually
discarded become a \emph{spatio-temporal filter tool}. The LLM's job is
not to look at pixels: it is to decide \emph{which tool to call next}
given the current evidence, and to declare \emph{when} the answer is
trustworthy.

Building on this insight we propose \method, the first LLM agent for
multi-camera person search. \method{} is built around a Planner /
Executor / Critic loop driven by a frontier LLM. The Planner emits one
function call at each step, choosing among $10$ tools that wrap the gait
encoder, the CLIP-Re-ID encoder, a multi-vector reranker, a vision-language
captioner, a spatio-temporal filter, a quality scorer, a FAISS coarse
retriever, and an evidence fuser; the Executor runs the tool and returns a
typed result; the Critic scores the step's usefulness and the running
confidence, triggering a \emph{self-correcting} replan whenever
modalities disagree or confidence stagnates. We evaluate
\method{} on MEVID~\cite{davila2023mevid}, MARS~\cite{zheng2016mars},
DukeMTMC-VideoReID~\cite{wu2018duke}, and the clothing-change benchmarks
PRCC~\cite{yang2019prcc} and LTCC~\cite{qian2020ltcc}, using both standard
mAP/CMC and three new metrics tailored to agent-based retrieval:
\emph{Reasoning Faithfulness}, \emph{Tool-call Efficiency}, and
\emph{Cross-camera Recall}. To summarize, our contributions are:

\begin{itemize}\setlength\itemsep{2pt}
  \item \textbf{Task reformulation.} We are the first to cast multi-camera
        person search as a \emph{tool-using LLM agent task}, replacing the
        prevailing single-shot matching paradigm with an explicit
        plan--act--criticize loop.
  \item \textbf{Tool library and schema.} We design a $10$-tool library
        (gait, Re-ID, VLM caption, spatio-temporal filter, quality scorer,
        FAISS, multi-vector reranker, evidence fuser, SAM silhouette,
        track extractor) with strict JSON schemas, so the Planner can
        compose them safely without function-calling support from the
        underlying LLM gateway.
  \item \textbf{Spatio-temporal reasoning prompt.} We integrate the
        camera-topology graph and time-window prior into the Planner
        prompt, so the agent prunes physically infeasible candidates
        before invoking heavy encoders.
  \item \textbf{Self-correcting agent loop.} We introduce a Critic that
        outputs a calibrated confidence and a \emph{needs-correction}
        flag, enabling the system to expand cameras, broaden time
        windows, or switch encoders when evidence is weak. We further
        prevent pathological loops with a tool-argument deduplication
        guard.
  \item \textbf{New evaluation protocol.} Beyond the standard mAP/CMC, we
        propose \emph{Reasoning Faithfulness} (alignment between stated
        rationale and actual tool output) and \emph{Tool-call Efficiency}
        (average number of calls per correctly answered query), turning
        the otherwise opaque agent behaviour into a measurable artefact.
\end{itemize}

Empirically, \method{} achieves \todo{X.X}\% mAP on MEVID, an absolute
improvement of \todo{$+\Delta$}\% over the strongest hand-tuned
gait+Re-ID fusion baseline, while also raising rank-1 by \todo{X.X}\% on
PRCC and LTCC. The agent makes on average \todo{N} tool calls per query,
\todo{Y}\% of which are validated by the Critic as useful; the resulting
reasoning chains are rated as faithful by \todo{Z}\% of human evaluators.
The remainder of the paper is organized as follows.
Section~\ref{sec:related} reviews video Re-ID, multi-camera person search,
and tool-using LLM agents. Section~\ref{sec:method} describes the tool
library, the Planner / Critic prompts, and the self-correcting loop.
Section~\ref{sec:exp} reports our experiments and ablations. We conclude
in Section~\ref{sec:conclusion}.
